\subsection{Message delays}

\textbf{Setup.}
We delayed the \texttt{promise} and/or \texttt{vote} replies by $D$~ms in the
acceptor using \texttt{timer:send\_after/3}. The state is updated when the
request arrives, so the delay models network latency. The protocol was left
unchanged: 5 acceptors, 3 proposers (sleeps 0/100/200~ms), timeout 2000~ms,
backoff 10~ms. For automated runs the GUI panels were replaced by processes
that discard messages; each point is the mean of five runs (ten around the
timeout), each in a fresh Erlang VM. A trial without a decision after 12~s is
a failure.

%%TABLE_MAIN%%

\textbf{Results.}
Consensus time grows linearly with $D$ (Table~\ref{tab:delays}). Delaying only
votes adds about $D$: promises are immediate, so the first proposer wins.
Delaying only promises adds about $D+200$~ms, because the later proposers
(higher sequence numbers) preempt the first one while its promises are in
flight; the last proposer wins and the extra 200~ms is its initial sleep.
Delaying both pays the delay twice, about $2D+200$~ms. Agreement was never
violated: the winner always decided in its first ballot, while losing
proposers needed two or three.

At $D=2000$~ms and above no run decided. Replies miss the timeout, every phase
aborts, and retrying does not help, so delays affect liveness but not safety.

\textbf{Timeout.}
Termination requires $D<\mathrm{timeout}$ \emph{per phase}, not
$2D<\mathrm{timeout}$: with $D=1500$ in both directions the round trip is
3~s, but a 2000~ms timeout still succeeds because each phase has its own
budget (Table~\ref{tab:timeout}). At $D=2500$ a 4000~ms timeout is needed.
Increasing the timeout does not reduce the consensus time, but it slows down
failure detection and makes losing proposers wait longer. Since delays are
unbounded, no finite timeout guarantees termination.

%%TABLE_TIMEOUT%%
